As a first experiment, we use the boolean model on the full Broad dataset,
binarized in pre-processing via quantile thresholding.
The idea is to allow a moderate ruleset complexity (meaning a reasonably high $C$)
and study the performance of the ruleset in terms of accuracy.
At this point we are not interested in preventing overfitting, but merely the existence of a performant ruleset.

Let $\bm{v}$ be a vector containing measurements.
Define $Q_{\bm{v}}\colon[0,1] \to \R$ as the quantile function (or percentile function) over the measurements in $\bm{v}$.
The quantile function uses linear interpolation when given an unobserved quantile level as input value.
The target boolean for point $i$ is
\[y_i \geq Q_{\bm{y}}(Q_y),\]
where $Q_y \in [0,1]$ is the chosen target quantile level.

Note that the choice of $Q_y$ determines the balance of positive and negative samples in the dataset.
For $Q_y = 0.5,$ the target is simply whether a point's PD-L1 expression is at least the median.
By construction, this leads to a (nearly) perfectly balanced dataset.
For $Q_y = 0.9$, on the other hand, the dataset is heavily imbalanced, as only a tenth of the samples are positive.
It is important to keep this in mind when evaluating models.
Notably, the empty ruleset $\emptyset$ achieves an accuracy of $Q_y$ for all $Q_y \in [0,1]$.
This is because the empty ruleset covers no points and a fraction of $Q_y$ points is negative by construction.
Hence, a ruleset for $Q_y = 0.9$, should show at least $0.9$ accuracy, and ideally much higher.
Other performance measures that may be considered in the case of imbalanced datasets are precision and recall.
Precision measures the fraction of positively classified points that is truly positive.
Recall measures the fraction of positive points that was positively classified.
%\begin{definition}[Precision]
%    The precision of a ruleset $S$ is
%    \[
%        \text{Precision}(S) = \frac
%        {\left|\left\{i \in \I : \hat{y}_S(\bm{X}_i) = 1 \wedge y_i = 1\right\}\right|}
%        {\left|\left\{i \in \I : \hat{y}_S(\bm{X}_i) = 1\right\}\right|}
%    \]
%\end{definition}

Denote by $\nq$ the number of thresholds to use for each TF\@.
Then for some TF with observed activity score vector $\bm{a},$ we consider the binary features
\begin{equation}\label{eq:BM_TF_activity_quantiles}
    \begin{split}
        a_{\cdot} &\geq Q_{\mathbf{a}}\left(\frac{1}{n_\text{quantiles}+1}\right), \\
        a_{\cdot} &\geq Q_{\mathbf{a}}\left(\frac{2}{n_\text{quantiles}+1}\right), \\
        \vdots & \tquad \vdots \\
        a_{\cdot} &\geq Q_{\mathbf{a}}\left(\frac{n_\text{quantiles}}{n_\text{quantiles}+1}\right).
    \end{split}
\end{equation}
Here, $\cdot$ is a placeholder for a point $i$.
All these binary features combined form $\bm{X}.$

We consider two variants of the experiment.
In the first variant, only the binary features described in \eqref{eq:BM_TF_activity_quantiles} are used.
This means that rules with more than one TF represent simultaneous activity of these TF~.
The second variant also adds the reverse binary features, corresponding to the upper-bounds
\begin{equation}\label{eq:BM_TF_activity_quantiles_reverse}
    \begin{split}
        a_{\cdot} &\leq Q_{\mathbf{a}}\left(\frac{1}{n_\text{quantiles}+1}\right), \\
        a_{\cdot} &\leq Q_{\mathbf{a}}\left(\frac{2}{n_\text{quantiles}+1}\right), \\
        \vdots & \tquad \vdots \\
        a_{\cdot} &\leq Q_{\mathbf{a}}\left(\frac{n_\text{quantiles}}{n_\text{quantiles}+1}\right).
    \end{split}
\end{equation}
A rule can then be a combination of lower- and upper-bounds.
Assume for now that a TF appears at most once in a rule.
Then a rule can represent simultaneous activity of some TF and inactivity of others.
An upper-bound can be interpreted as an ``interfering'' TF that, when present,
prevents other TF from acting together.

We solve the model (via the price-then-branch heuristic) for
\[C = 20, \quad D = 3, \quad \nq = 1, 3, 5, 7, 9. \]
where $C$ is the maximum ruleset complexity and $D$ is the maximum rule size.
Recall that, in accordance with~\cite{lawless_interpretable_2023}, a rule's complexity is $1$ plus the number of features in it.
As such, a ruleset complexity of $20$ is equivalent to five rules of size three.
A ruleset complexity of $18$ is equivalent to six rules of size two.

We employ a heuristic pricer with beam widths $(1000, 500, 500)$.
We also employ the exact pricer.
We limit pricing (the finding of rules) to 23 hours and branching to 1 hour.
